\documentclass[11pt,a4paper]{article}

\usepackage[T1]{fontenc}
\usepackage[utf8]{inputenc}
\usepackage{lmodern}
\usepackage[margin=2.3cm]{geometry}
\usepackage{xcolor}
\usepackage{booktabs}
\usepackage{tabularx}
\usepackage{enumitem}
\usepackage{listings}
\usepackage[most]{tcolorbox}
\usepackage{amssymb}
\usepackage[colorlinks=true,linkcolor=blue!50!black,urlcolor=blue!50!black]{hyperref}

% ---------------------------------------------------------------------------
% Edit these before sending the document out.
\newcommand{\coursedate}{\textbf{[DATE]}}
\newcommand{\instructorcontact}{\textbf{[NAME, EMAIL]}}
\newcommand{\courserepo}{\detokenize{git@github.com:marcelojbp/physics_RSE_and_AI.git}}  % SSH address
% ---------------------------------------------------------------------------

\definecolor{codebg}{gray}{0.95}
\lstset{
  basicstyle=\ttfamily\small,
  backgroundcolor=\color{codebg},
  frame=single, rulecolor=\color{gray!40},
  breaklines=true, columns=fullflexible, keepspaces=true,
  showstringspaces=false, upquote=true, escapeinside={(*}{*)},
  xleftmargin=3pt, xrightmargin=3pt, aboveskip=6pt, belowskip=6pt
}

\newtcolorbox{note}[1][Note]{colback=blue!4, colframe=blue!45!black,
  fonttitle=\bfseries, title=#1, boxrule=0.6pt, arc=2pt}
\newtcolorbox{warn}[1][Watch out]{colback=orange!6, colframe=orange!70!black,
  fonttitle=\bfseries, title=#1, boxrule=0.6pt, arc=2pt}

\newcommand{\cmd}[1]{\texttt{#1}}

\setlist{itemsep=2pt, topsep=4pt}
\setlength{\parindent}{0pt}
\setlength{\parskip}{5pt}

\title{\textbf{AI for Physics Projects}\\[4pt]
\large Setup instructions for macOS, Windows and Linux}
\author{}
\date{Please complete \emph{before} the course on \coursedate}

\begin{document}
\maketitle
\vspace{-1.5em}

\begin{note}[Why this matters]
The hands-on session lasts one hour. If we spend it installing software, we lose it.
Plan for \textbf{45--60 minutes}, most of it unattended while Julia packages
download and compile. If anything fails, send the error message to
\instructorcontact{} before the course; do not wait until the day.
\end{note}

\tableofcontents

% ===========================================================================
\section{What you will install, and why}

\begin{center}
\renewcommand{\arraystretch}{1.25}
\begin{tabularx}{\textwidth}{@{}l l X@{}}
\toprule
\textbf{Tool} & \textbf{Required?} & \textbf{Role in the course} \\
\midrule
Git              & required & Version control: the undo button and lab notebook that make it safe to let an agent edit your files. \\
GitHub account   & required & Hosts repositories; runs automated tests (continuous integration, CI). \\
GitHub CLI (\cmd{gh}) & required & Logs you into GitHub and sets up an SSH key for Git in one step. \\
Julia (via \cmd{juliaup}) & required & The language of all examples. \\
Julia packages   & required & DrWatson (project structure), QuantumOptics, OrdinaryDiffEq, LsqFit, CairoMakie, JLD2, Revise. \\
VS Code + Julia extension & required & Editor; also the easiest way to \emph{read diffs}, a central skill of the course. \\
AI coding agent  & required & An LLM tool that works directly in a folder on your computer (e.g.\ Claude Code or Codex). See Section~\ref{sec:agent}. \\
Course repository & required & Slides, exercises and the hands-on project; you clone it at the end of the setup. \\
GitHub Desktop   & optional & Graphical Git client (macOS and Windows only). \\
\bottomrule
\end{tabularx}
\end{center}

\textbf{Requirements.}
\begin{itemize}
  \item About \textbf{6\,GB} of free disk space (Julia packages are the bulk of it).
  \item Permission to install software. On a university-managed laptop you may need
        IT support; ask \emph{now}, not the day before.
  \item A working internet connection without a restrictive proxy. If your institute
        uses one, see Section~\ref{sec:troubleshooting}.
\end{itemize}

\textbf{Order.} Do Section~\ref{sec:accounts} (accounts), then \emph{only} the
section for your operating system (\ref{sec:macos}, \ref{sec:windows} or
\ref{sec:linux}), then Sections~\ref{sec:config}--\ref{sec:verify}, which are the
same for everyone.

\textbf{Terminal.} Several steps use a terminal:
\begin{itemize}
  \item \textbf{macOS:} \emph{Terminal} (Applications $\to$ Utilities).
  \item \textbf{Windows:} \emph{Terminal} or \emph{PowerShell} from the Start menu.
        Do \emph{not} use the old \emph{Command Prompt} (\cmd{cmd.exe}).
  \item \textbf{Linux:} your usual terminal emulator.
\end{itemize}
Type or paste each command, then press Enter. Lines starting with \cmd{\#} are
comments; you do not need to type them.

% ===========================================================================
\section{Accounts}\label{sec:accounts}

\subsection{GitHub}
\begin{enumerate}
  \item Create a free account at \url{https://github.com/signup}.
        Use an email address you will keep; your university address also qualifies
        you for free extras via \url{https://education.github.com}.
  \item Note your \textbf{username} and the \textbf{email} you registered with;
        you will need both in Section~\ref{sec:gitconfig}.
\end{enumerate}

\subsection{An AI coding agent}\label{sec:agent}
You need access to an LLM with \emph{local agent} capabilities: a tool that runs on
your computer, inside a project folder, and can read and edit files and run commands
(with your permission). A chat window in the browser is \emph{not} enough.
Examples are \textbf{Claude Code} (Anthropic) and \textbf{Codex} (OpenAI); use
whichever matches the subscription you already have.

Install it following its own instructions, and make sure that before the course you
can (i) start it in a folder and (ii) log in. Section~\ref{sec:agent-test} has a
one-minute test. The course does not depend on which agent you use.

% ===========================================================================
\section{macOS}\label{sec:macos}

Works on Apple Silicon (M1 and later) and Intel Macs, macOS 13 or newer.

\subsection{Homebrew (package manager)}
Homebrew installs and updates developer tools with one command each. Check
whether you already have it:
\begin{lstlisting}
brew --version
\end{lstlisting}
If this prints a version, skip to Section~\ref{sec:mac-tools}. Otherwise install it:
\begin{lstlisting}
/bin/bash -c "$(curl -fsSL https://raw.githubusercontent.com/Homebrew/install/HEAD/install.sh)"
\end{lstlisting}
It will ask for your Mac password (nothing appears as you type; that is normal) and
may install the \emph{Xcode Command Line Tools}, which takes a few minutes.

\begin{warn}[Apple Silicon: read the last lines of the output]
At the end, Homebrew prints ``Next steps'' with two or three commands that add
\cmd{brew} to your \cmd{PATH} (they contain \cmd{/opt/homebrew/bin/brew shellenv}).
\textbf{Copy and run them exactly as printed}, then close and reopen Terminal.
\end{warn}

\subsection{Git, GitHub CLI, VS Code, GitHub Desktop}\label{sec:mac-tools}
\begin{lstlisting}
brew install git gh
brew install --cask visual-studio-code
brew install --cask github          # optional: GitHub Desktop
\end{lstlisting}
macOS ships an older Git with the Xcode tools; the Homebrew version takes precedence
once \cmd{brew} is on your \cmd{PATH}.

\subsection{Julia}
Install \cmd{juliaup}, the official Julia version manager:
\begin{lstlisting}
curl -fsSL https://install.julialang.org | sh
\end{lstlisting}
Accept the defaults (press Enter). \textbf{Close and reopen Terminal}, then check:
\begin{lstlisting}
julia --version
\end{lstlisting}
Do not install Julia through \cmd{brew install julia} or from a \cmd{.dmg}; use
\cmd{juliaup} so everyone has the same setup. (\cmd{brew install juliaup} is also fine,
but install \cmd{juliaup} only one way, not both.)

Continue with Section~\ref{sec:config}.

% ===========================================================================
\section{Windows}\label{sec:windows}

Windows 10 (version 1809 or later) or Windows 11. All commands below run in
\textbf{PowerShell} or \textbf{Windows Terminal}, not in Command Prompt.
We use \cmd{winget}, the package manager built into Windows. Check it works:
\begin{lstlisting}
winget --version
\end{lstlisting}
If the command is not found, install \emph{App Installer} from the Microsoft Store
and reopen the terminal.

\subsection{Git}
\begin{lstlisting}
winget install --id Git.Git -e --source winget
\end{lstlisting}
This installs \emph{Git for Windows}, which includes \emph{Git Bash}. Several AI agents
use Git Bash on Windows, so do not skip this step even if you plan to use GitHub Desktop.
If you install Git from \url{https://git-scm.com} instead, the default options are
fine; in particular keep \emph{``Git from the command line and also from 3rd-party
software''} and \emph{``Checkout Windows-style, commit Unix-style line endings''}.

\subsection{GitHub CLI, VS Code, GitHub Desktop}
\begin{lstlisting}
winget install --id GitHub.cli -e
winget install --id Microsoft.VisualStudioCode -e
winget install --id GitHub.GitHubDesktop -e     # optional
\end{lstlisting}

\subsection{Julia}
\begin{lstlisting}
winget install --name Julia --id 9NJNWW8PVKMN -e -s msstore
\end{lstlisting}
This installs \cmd{juliaup} together with the current stable Julia. If the Microsoft
Store is blocked on your machine, download the installer from
\url{https://julialang.org/downloads/} instead.

\begin{warn}[Reopen the terminal]
After installing, \textbf{close all terminal windows and open a new one}; otherwise the
new commands are not found. Then check:
\begin{lstlisting}
git --version
gh --version
code --version
julia --version
\end{lstlisting}
\end{warn}

Continue with Section~\ref{sec:config}.

% ===========================================================================
\section{Linux}\label{sec:linux}

Commands are given for Ubuntu/Debian, Fedora and Arch. Use the line for your
distribution.

\subsection{Git and curl}
\begin{lstlisting}
sudo apt update && sudo apt install -y git curl     # Ubuntu / Debian
sudo dnf install -y git curl                        # Fedora
sudo pacman -S --needed git curl                    # Arch
\end{lstlisting}

\subsection{GitHub CLI}
\begin{lstlisting}
sudo dnf install -y gh                              # Fedora
sudo pacman -S github-cli                           # Arch
\end{lstlisting}
On \textbf{Ubuntu/Debian}, the version in the default repositories is often outdated.
Follow the official instructions (one copy-paste block) at
\url{https://github.com/cli/cli/blob/trunk/docs/install_linux.md}.

\subsection{VS Code}
\begin{itemize}
  \item \textbf{Ubuntu/Debian:} download the \cmd{.deb} file from
        \url{https://code.visualstudio.com/download}, then
\begin{lstlisting}
sudo apt install ./code_*.deb
\end{lstlisting}
  \item \textbf{Fedora:} download the \cmd{.rpm} file from the same page, then
        \cmd{sudo dnf install ./code-*.rpm}.
  \item \textbf{Arch:} \cmd{visual-studio-code-bin} from the AUR.
\end{itemize}
These packages add Microsoft's repository, so VS Code updates with your system.
GitHub Desktop has no official Linux version; use VS Code's built-in Git view instead.

\subsection{Julia}
\begin{lstlisting}
curl -fsSL https://install.julialang.org | sh
\end{lstlisting}
Accept the defaults, \textbf{close and reopen the terminal}, and check
\cmd{julia --version}. Do not use the \cmd{julia} package from your distribution;
it is usually outdated or built differently.

Continue with Section~\ref{sec:config}.

% ===========================================================================
\section{Configuration (all systems)}\label{sec:config}

\subsection{Git}\label{sec:gitconfig}
Tell Git who you are. Use the \textbf{same email as your GitHub account}, otherwise
your commits will not be linked to you.
\begin{lstlisting}
git config --global user.name "Your Name"
git config --global user.email "you@example.com"
git config --global init.defaultBranch main
git config --global core.editor "code --wait"
git config --global pull.rebase false
\end{lstlisting}
The last three are recommended defaults: new repositories start on a branch called
\cmd{main}, matching GitHub; when Git needs you to write a message it opens VS Code
instead of \cmd{vim}; and \cmd{git pull} merges instead of stopping with
``\emph{Need to specify how to reconcile divergent branches}''.
Line endings (one command, depending on your system):
\begin{lstlisting}
git config --global core.autocrlf input     # macOS and Linux
git config --global core.autocrlf true      # Windows
\end{lstlisting}
Check your settings with \cmd{git config --global --list}.

\subsection{Connect Git to GitHub (SSH)}\label{sec:ssh}
In this course Git talks to GitHub over \textbf{SSH}: your computer holds a private
key, GitHub holds the matching public key, and no password is ever typed. The GitHub
CLI creates the key and registers it with GitHub in one go:
\begin{lstlisting}
gh auth login
\end{lstlisting}
Answer the questions as follows:
\begin{enumerate}
  \item Where do you use GitHub? $\to$ \textbf{GitHub.com}
  \item Preferred protocol for Git operations? $\to$ \textbf{SSH}
  \item Generate a new SSH key to add to your GitHub account? $\to$ \textbf{Yes}\\
        (If you already have a key, it asks instead whether to \emph{upload} your
        existing public key; choose that key.)
  \item Enter a passphrase for your new SSH key $\to$ press \textbf{Enter} to leave it
        empty, unless you already know how to use an SSH agent (see the box below).
  \item Title for your SSH key $\to$ accept the default, or type the name of this
        computer.
  \item How would you like to authenticate? $\to$ \textbf{Login with a web browser}
\end{enumerate}
Copy the one-time code it shows, press Enter, paste the code in the browser and
approve. Then check both the GitHub login and the SSH connection:
\begin{lstlisting}
gh auth status
ssh -T git@github.com
\end{lstlisting}
\cmd{gh auth status} should say \emph{``Logged in to github.com account \ldots''} and
\emph{``Git operations protocol: ssh''}. The first time you run \cmd{ssh -T}, it asks
\emph{``Are you sure you want to continue connecting?''}: type \textbf{yes}. It should
then answer \emph{``Hi \textless username\textgreater! You've successfully
authenticated, but GitHub does not provide shell access.''} That message means it
worked.

\begin{note}[Passphrase or not?]
The key file is what proves you are you, so anyone who copies it can push as you.
A passphrase protects it, but you then have to unlock it once per session.
On \textbf{macOS} this is painless: add these lines to the file \cmd{\textasciitilde/.ssh/config}
(create it if it does not exist) and macOS remembers the passphrase in its keychain:
\begin{lstlisting}
Host github.com
  AddKeysToAgent yes
  UseKeychain yes
\end{lstlisting}
On Windows and Linux, a passphrase requires running an SSH agent; if you have never
done this, leave the passphrase empty for now.
\end{note}

\textbf{Cloning.} Always copy the \textbf{SSH} address of a repository
(\cmd{git@github.com:user/repo.git}), not the HTTPS one: on GitHub, click the green
\emph{Code} button and choose the \emph{SSH} tab. \cmd{gh repo clone user/repo} uses
SSH automatically.

\subsection{VS Code}
\begin{enumerate}
  \item \textbf{macOS only:} open VS Code, press \cmd{Cmd+Shift+P}, type
        \emph{shell command}, and choose
        \emph{``Shell Command: Install `code' command in PATH''}.
        (Windows and Linux do this automatically.)
  \item Install the Julia extension from a terminal:
\begin{lstlisting}
code --install-extension julialang.language-julia
\end{lstlisting}
        Many AI agents also offer a VS Code extension; installing it is optional.
  \item Open VS Code once and click the \emph{Source Control} icon in the left bar.
        This is where we will read diffs during the course.
\end{enumerate}

\subsection{Course material and Julia packages}\label{sec:julia-pkgs}
This step copies the course repository to your computer (this also tests your SSH
connection), then downloads and compiles everything the hands-on session needs. It
takes \textbf{10--20 minutes}; start it and let it run. Create a folder for the
course in your home folder and clone the repository into it:
\begin{lstlisting}
cd ~
mkdir ai-course
cd ai-course
git clone (*\courserepo*) course
cd course
julia setup/setup_julia.jl
\end{lstlisting}
The script ends with \emph{``All good: Julia \ldots{} is ready for the course.''}
It installs the packages in the table below, then the exact package versions of the
hands-on project \cmd{DrivenSpin}, and runs that project's tests.

Keep the folder \cmd{ai-course}: we start from it on the day, and update the course
material with \cmd{git pull} inside \cmd{course}. Do not edit files inside
\cmd{course}; in the hands-on you work on your own copy.

\begin{warn}[Keep this Julia version until the course]
Packages are compiled for one Julia version. If you update Julia (e.g.\ with
\cmd{juliaup update}) after this step, everything is compiled again on the day.
\end{warn}

\textbf{If cloning fails} (and Section~\ref{sec:troubleshooting} does not help), still
prepare the packages: start Julia by typing \cmd{julia} in a terminal, then paste the
following at the \cmd{julia>} prompt, and send the clone error to \instructorcontact.
\begin{lstlisting}[language={}]
using Pkg
Pkg.add(["DrWatson", "Revise"])
Pkg.activate(temp=true)
Pkg.add(["DrWatson", "QuantumOptics", "OrdinaryDiffEq",
         "LsqFit", "CairoMakie", "JLD2"])
Pkg.precompile()
using DrWatson, QuantumOptics, OrdinaryDiffEq, LsqFit, CairoMakie, JLD2
println("All packages loaded.")
\end{lstlisting}
Leave Julia with \cmd{exit()}.

What the packages are for:
\begin{center}
\renewcommand{\arraystretch}{1.15}
\begin{tabularx}{\textwidth}{@{}l X@{}}
\toprule
DrWatson        & Project structure, reproducible file names, saving results tagged with the Git commit. \\
QuantumOptics   & Reference solver for spins and cavities; used as an independent check. \\
OrdinaryDiffEq  & Time evolution (the ODE part of DifferentialEquations.jl, lighter to install). \\
LsqFit          & Fitting synthetic ``experimental'' data. \\
CairoMakie      & Plotting. \\
JLD2            & File format DrWatson uses to save results. \\
Revise          & Reloads your code automatically as you (or the agent) edit it. \\
\bottomrule
\end{tabularx}
\end{center}
\cmd{Test} and \cmd{LinearAlgebra} ship with Julia and need no installation.

DrWatson and Revise go into your default environment, so they are always available.
The other packages are downloaded into Julia's shared cache; your copy of the hands-on
project reuses them, so you will not wait for compilation on the day.

\subsection{AI agent: one-minute test}\label{sec:agent-test}
Agents work inside a folder, so create an empty practice folder and start your agent
there (from the terminal or its app, depending on the tool):
\begin{lstlisting}
mkdir ai-course-test
cd ai-course-test
\end{lstlisting}
Log in if asked, then give it a first request:
\begin{lstlisting}[language={}]
Write a Julia script hello.jl that prints the Julia version, and run it.
\end{lstlisting}
The test passes if the agent creates \cmd{hello.jl} in that folder, runs it with
\emph{your} Julia, and shows the version printed by \cmd{julia --version}. It will
usually ask permission before creating the file and before running it; approve both.
In the course the agent also runs Git, so give it a second request:
\begin{lstlisting}[language={}]
Make this folder a git repository and commit hello.jl. Then show me git log.
\end{lstlisting}
The log should show one commit with your name. You can delete the folder afterwards.

% ===========================================================================
\section{Final check}\label{sec:verify}

In a \textbf{new} terminal, run each of these. Every line should print a version or
``logged in'' and no error.
\begin{lstlisting}
git --version
git config user.email
gh auth status
ssh -T git@github.com
code --version
julia --version
julia -e "using DrWatson"
cd ~/ai-course/course
git status
\end{lstlisting}
\cmd{ssh -T} should greet you by your GitHub username (see Section~\ref{sec:ssh}).
\cmd{julia -e "using DrWatson"} prints nothing when it works and an error when it does
not. \cmd{git status} should say \emph{``On branch main. Your branch is up to date with
`origin/main'.''}
Then tick the list:

\begin{center}
\renewcommand{\arraystretch}{1.3}
\begin{tabularx}{\textwidth}{@{}c X@{}}
$\square$ & I have a GitHub account and \cmd{gh auth status} says I am logged in, with protocol \cmd{ssh}. \\
$\square$ & \cmd{ssh -T git@github.com} answers ``Hi \textless username\textgreater! You've successfully authenticated\ldots''. \\
$\square$ & \cmd{git config user.email} shows the email of my GitHub account. \\
$\square$ & VS Code opens and has the Julia extension. \\
$\square$ & The course repository is in \cmd{ai-course/course}, and \cmd{julia setup/setup\_julia.jl} finished with ``All good''. \\
$\square$ & My AI agent starts in a folder, I am logged in, and it created, ran and committed \cmd{hello.jl}. \\
\end{tabularx}
\end{center}

Bring your laptop \textbf{and its charger} on \coursedate.

% ===========================================================================
\section{Troubleshooting}\label{sec:troubleshooting}

\begin{description}[style=nextline, leftmargin=0pt, itemsep=6pt]
\item[``command not found'' / ``is not recognized'' right after installing]
  Close \emph{all} terminal windows and open a new one. On Windows, if that fails,
  sign out and back in. On macOS, check that the Homebrew ``Next steps'' commands were run.

\item[Julia package installation stalls or fails with network errors]
  Usually a proxy or firewall. Try from home or eduroam. If your institute requires a
  proxy, set it before starting Julia (ask IT for the address):
\begin{lstlisting}
export HTTPS_PROXY=http://proxy.example.org:8080     # macOS / Linux
$env:HTTPS_PROXY="http://proxy.example.org:8080"     # Windows PowerShell
\end{lstlisting}
  Then run the installation again; it resumes where it stopped.

\item[Precompilation is very slow]
  Expected on the first run, especially for CairoMakie and OrdinaryDiffEq. Keep the
  laptop plugged in and awake. It only happens once.

\item[The agent cannot find \cmd{julia} or \cmd{git}]
  Start the agent from a \emph{new} terminal opened after installing everything, so it
  inherits the updated \cmd{PATH}.

\item[\cmd{ssh -T git@github.com} says ``Permission denied (publickey)'']
  Your key is missing or not registered with GitHub. Run
  \cmd{gh ssh-key list} to see the keys GitHub knows about. If the list is empty, run
  \cmd{gh auth refresh -s admin:public\_key} and then
  \cmd{gh ssh-key add \textasciitilde/.ssh/id\_ed25519.pub}.

\item[\cmd{ssh -T} hangs or says ``Connection timed out'']
  Many institute networks block SSH's usual port (22). GitHub also accepts SSH on
  port 443. Add these lines to \cmd{\textasciitilde/.ssh/config} (create the file if needed;
  on Windows it is \cmd{.ssh\textbackslash config} in your user folder):
\begin{lstlisting}
Host github.com
  Hostname ssh.github.com
  Port 443
  User git
\end{lstlisting}
  and run \cmd{ssh -T git@github.com} again. If you already have a \cmd{Host github.com}
  block (from the passphrase box), add the three lines to it.

\item[\cmd{git clone} says ``Repository not found'' or ``Could not read from remote repository'']
  Check the SSH connection first with \cmd{ssh -T git@github.com} (see the two items
  above). If that works, check that you copied the address exactly.

\item[\cmd{git push} asks for a username or password]
  The repository was cloned with its HTTPS address. Switch it to SSH, from inside the
  repository folder:
\begin{lstlisting}
git remote set-url origin git@github.com:<user>/<repo>.git
\end{lstlisting}
  and make SSH the default for the GitHub CLI: \cmd{gh config set git\_protocol ssh}.

\item[University-managed laptop, no admin rights]
  \cmd{juliaup} installs into your user folder and does not need admin rights. Git, VS Code and the GitHub CLI usually do; ask IT, or on Windows use
  \cmd{winget install} with \cmd{--scope user} where available.
\end{description}

If you are stuck, send the \emph{exact} error message (a screenshot is fine) and your
operating system to \instructorcontact.

\end{document}
